%Chargement des paquets
\usepackage{amsmath}
\usepackage{amsfonts}
\usepackage{amssymb}
\usepackage{enumerate}
\usepackage{mathtools}
\usepackage{bbm}
\usepackage{xparse, etoolbox}
\usepackage{enumerate}
\usepackage{enumitem}
\usepackage{mathabx}
\usepackage{babel}[french]

%environment
\usepackage{amsthm}
\usepackage{keytheorems}
\usepackage{hyperref} 

%Interface théorème
\newkeytheorem{lem}[
	name = Lemme
]

\newkeytheorem{prop}[
	name = Proposition
]

\newkeytheorem{thm}[
	name = Théorème
]

\newkeytheorem{coro}[
	name = Corollaire
]

\renewcommand*{\proofname}{Démonstration}

\theoremstyle{definition}
\newtheorem{defn}{Définition}

\theoremstyle{definition}
\newtheorem*{nota}{Notation}

\theoremstyle{definition}
\newtheorem*{limi}{Liminaire}

\theoremstyle{remark}
\newtheorem{rmq}{Remarque}

\theoremstyle{plain}
\newtheorem*{fait}{Fait}


%Ensembles de nombres
\newcommand{\N} {\mathbb{N}}
\newcommand{\Ne}{\N^\ast}
\newcommand{\Z} {\mathbb{Z}}
\newcommand{\Q} {\mathbb{Q}}
\newcommand{\R} {\mathbb{R}}
\newcommand{\Rb}{\overline{\mathbb{R}}}
\newcommand{\Rp}{\R_+}
\newcommand{\Rm}{\R_-}
\newcommand{\K} {\mathbb{K}}
\newcommand{\Ket} {\mathbb{K^{\ast}}}
\newcommand{\Cx}{\mathbb{C}}

%Ensembles, tribus
\newcommand{\A}{\mathbb{A}}
\renewcommand{\P}{\mathcal{P}}
\newcommand{\T}{\mathcal{T}}
\newcommand{\F}{\mathcal{F}}
\newcommand{\C} {\mathcal{C}}
\newcommand{\ouv}{\mathcal{O}}
\newcommand{\B}{\mathcal{B}}
\newcommand{\M}{\mathcal{M}}
\newcommand{\etag}{\mathcal{E}}
\newcommand{\etagOp}{\etag^{+}(\Omega)}
\newcommand{\etagO}{\etag(\Omega)}
%%Lp avant quotientage
\newcommand{\Lic}{\mathcal{L}^1}
\newcommand{\Lpc}{\mathcal{L}^p}
\newcommand{\Linfc}{\mathcal{L}^{\infty}}
\newcommand{\Lifull}{\Lic(\Omega, \B, m)}
\newcommand{\Lpfull}{\Lpc(\Omega, \B, m)}
\newcommand{\Linffull}{\Linfc(\Omega, \B, m)}
%%Lp après quotientage
\newcommand{\Li}{L^1}
\newcommand{\Ld}{L^2}
\newcommand{\Lp}{L^p}
\newcommand{\Linf}{L^{\infty}}
\newcommand{\Listd}{\Li(\Omega, m)} 
\newcommand{\Ldstd}{\Ld(\Omega, m)} 
\newcommand{\Lpstd}{\Lp(\Omega, m)} 

%Quelques opérateurs
\DeclareMathOperator{\codim}{codim}
\DeclareMathOperator{\rg}{rg}
\DeclareMathOperator{\tr}{tr}
\DeclareMathOperator{\vect}{Vect}
\DeclareMathOperator{\ima}{Im}
\DeclareMathOperator{\id}{Id}
\DeclareMathOperator{\sign}{sign}
\DeclareMathOperator{\ext}{ext}
\newcommand{\ind}{\mathbbm{1}}
\newcommand*{\spr}[2]{\left(\,#1\,\middle|\,#2\,\right)}
\newcommand{\interior}[1]{%
  {\kern0pt#1}^{\mathrm{o}}%
}
\DeclareMathOperator{\card}{card}
\DeclareMathOperator{\ppcm}{ppcm}

%Commandes de pure flemme
\newcommand{\veps}{\varepsilon}
\newcommand{\intab}{\int_{a}^{b}}
\newcommand{\intR}{\int_{-\infty}^{\infty}}
\newcommand{\intinf}{\int_{- \infty}^{+ \infty}}
\newcommand{\equi}{\Leftrightarrow}
\newcommand{\Kev}{$\K$-espace vectoriel }
\newcommand{\Kevs}{$\K$-espaces vectoriels }
\newcommand{\Rev}{$\R$-espace vectoriel }
\newcommand{\Revs}{$\R$-espaces vectoriels }
\newcommand{\Cev}{$\Cx$-espace vectoriel }
\newcommand{\Cevs}{$\Cx$-espaces vectoriels }
\newcommand{\evn}{(E, \norm{.})}
\newcommand{\interf}[2]{[#1,#2]}
\newcommand{\limb}{\lim\limits_}
\newcommand{\lebn}{\lambda_n}
\newcommand{\pp}{\text{~presque partout}}
\newcommand{\derivp}[2]{\frac{\partial #1}{\partial #2}}
\newcommand{\famiu}{(u_i)_{i \in I}}
\newcommand{\sev}{\text{~s.e.v.}}
\newcommand{\Mnp}{M_{n,p}(\K)}
\newcommand{\Mn}{M_{n}(\K)}
\newcommand{\tq}{\text{~t.q.~}}
\newcommand{\mq}{~montrer~que~}
\newcommand{\et}{\quad \text{et} \quad}
\newcommand{\ou}{\text{~ou~}}
\newcommand{\Kx}{\K[X]}


%Commandes de gars chelou
\renewcommand{\subset}{\subseteq}

%Commande restriction, ty duckduckgo
\def\restr#1#2{\mathchoice
              {\setbox1\hbox{${\displaystyle #1}_{\scriptstyle #2}$}
              \restraux{#1}{#2}}
              {\setbox1\hbox{${\textstyle #1}_{\scriptstyle #2}$}
              \restraux{#1}{#2}}
              {\setbox1\hbox{${\scriptstyle #1}_{\scriptscriptstyle #2}$}
              \restraux{#1}{#2}}
              {\setbox1\hbox{${\scriptscriptstyle #1}_{\scriptscriptstyle #2}$}
              \restraux{#1}{#2}}}
\def\restraux#1#2{{#1\,\smash{\vrule height .8\ht1 depth .85\dp1}}_{\,#2}} 

%Quelques normes
\DeclarePairedDelimiter\abs{\lvert}{\rvert}
\DeclarePairedDelimiter\labs{\bigl\lvert}{\bigr\rvert}
\DeclarePairedDelimiter\norm{\lVert}{\rVert}
\DeclarePairedDelimiterXPP\normi[1]{}\lVert\rVert{_1}{\ifblank{#1}{\:\cdot\:}{#1}}
\DeclarePairedDelimiterXPP\normd[1]{}\lVert\rVert{_2}{\ifblank{#1}{\:\cdot\:}{#1}}
\DeclarePairedDelimiterXPP\normp[1]{}\lVert\rVert{_p}{\ifblank{#1}{\:\cdot\:}{#1}}
\DeclarePairedDelimiterXPP\normq[1]{}\lVert\rVert{_q}{\ifblank{#1}{\:\cdot\:}{#1}}
\DeclarePairedDelimiterXPP\norminf[1]{}\lVert\rVert{_\infty}{\ifblank{#1}{\:\cdot\:}{#1}}

%Bracket en angle
\DeclarePairedDelimiter\angl{\langle}{\rangle}

%Set, parenthèses
\newcommand{\set}[1]{\{ #1 \}}
\newcommand{\bigset}[1]{\bigl\{ #1 \bigr\}}
\newcommand{\Bigset}[1]{\Bigl\{ #1 \Bigr\}}
\newcommand{\bigpar}[1]{\bigl( #1 \bigr)}
\newcommand{\Bigpar}[1]{\Bigl( #1 \Bigr)}

%Opitmisation des opérateurs de comparaison
\DeclarePairedDelimiter\tio{o (}{)}
\DeclarePairedDelimiter\gro{O (}{)}

\renewcommand{\L}{\mathbb{L}}
\renewcommand{\F}{\mathbb{F}}
\newcommand{\U}{\mathcal{U}}
\newcommand{\D}{\mathcal{D}}

\pagestyle{fancy}

\fancyhead[C]{}
\fancyhead[R]{}

\begin{document}

\underline{Source :} Rombaldi - Algèbre - pages 422, 425

\underline{Leçons :}
\begin{itemize}
\item  123 : Corps finis. Applications.
\item   141 : Polynômes irréductibles à une indéterminée. Corps de rupture. Exemples et applications.
\item   125 : Extensions de corps. Exemples et applications

\end{itemize}

\begin{nota}
Dans ce développement, $p \geq 2$ est un nombre premier.
On note $\U_n(p)$ l'ensemble des polynômes unitaires irréductibles de degré $n$ à coefficients dans $\F_p$
et $I_n(p)$ le cardinal de cet ensemble (éventuellement égal à 0).
Pour tout entier $n \in \Ne$, on note $P_n(X) = X^{p^n} - X \in \F_p[X]$.
\end{nota}

\begin{lem}
\label{lem:PdivPn}
Tout diviseur irréductible de $P_n$ dans $F_p[X]$ est de degré divisant $n$.
Réciproquement, pour tout diviseur $d$ de $n$, tout polynôme $P \in \U_d(p)$ divise le polynôme $P_n(X) = X^{p^n} - X$ 
dans $\F_p[X]$.
\end{lem}

\begin{proof}
Soit $P \in \U_d(p)$ (avec $d \leq n$) qui divise $P_n$.
On a $\overline{P_n} = \overline{0}$ dans le corps $\F^{p^d}=\F_p[X]/(P)$ et donc $\overline{X}^{p^n} = \overline{X}$. 
Comme $(\overline{X}^k)_{0 \leq k \leq d-1}$ est une base du $\F_p$-espace vectoriel $F_p[X]/(P)$, tout élément $Q$ de ce dernier
s'écrit sous la forme :
\[ Q = \sum_{k=0}^{d-1} a_k \overline{X}^k \qquad (a_k \in \F_p) , \]
et on déduit que :
\[ Q^{p^n} = \left ( \sum_{k=0}^{d-1} a_k \overline{X}^k \right )^{p^n} 
	= \sum_{k=0}^{d-1} a_k^{p^n} \left ( \overline{X}^k \right )^{p^n} , \]
car $p$ divise tous les coefficients binomiaux $\binom{p^n}{i}$ avec $1 \leq i \leq p^n-1$.
De plus $a_k^{p^n} = a_k$ dans $\F_p$ selon le théorème de Fermat 
et $\left ( \overline{X}^k \right )^{p^n} = \left ( \overline{X}^{p^n} \right)^k = \overline{X}^k$.
On conclut donc que :
\[ Q^{p^n} = \sum_{k=0}^{d-1} a_k \overline{X}^k = Q . \]
En conséquence, pour $Q \in \F_{p^d}^{\ast}$, on a $Q^{p^n-1} = 1$.
Le groupe multiplicatif $\F_{p^d}^{\ast}$ étant d'ordre $p^d-1$, pour tout $Q \in \F_{p^d}^{\ast}$ on a $Q^{p^d-1} = 1$.
Autrement dit, $p^d-1$ divise $p^n-1$ ce qui implique que $d$ divise $n$.
En effet, $n=dq+r$ avec $0 \leq r < d$ et donc :
\[ p^n-1 = p^{dq} p^r - 1 = (p^{dq}-1)p^r + p^r - 1 , \]
et le fait que $p^d-1$ divise $p^{dq}-1$ et $p^n-1$ entraîne qu'il divise $p^r-1$ soit $r=0$. \\

Soit $P \in \U_d(p), \F_{p^d} = \F_p[X]/(P)$ est un corps de cardinal $p^d$ (car $P$ est irréductible), 
donc $\F_{p^d}^{\ast}$ est un groupe multiplicatif de cardinal $p^d-1$,
et $\bar{X}^{p^d-1} = \bar{1}$ (théorème de Lagrange), donc $\bar{X}^{p^d} = \bar{X}$. Par récurrence, on peut montrer que 
$\bar{X}^{p^{kd}} = \bar{X}$ pour tout $k \in \N$ (en effet $\bar{X}^{p^{(k+1)d}} = \bigpar { \bar{X}^{p^{kd}} }^{p^d}$). \\
Comme $d$ divise $n$, il existe $q$ tel que $n=dq$ et par conséquence, $\bar{X}^{p^n} = \bar{X}^{p^{qd}} = \bar{X}$, 
autrement dit, $\widebar{P_n} = \bar{0}$ dans $F_{p^d}$, ce qui démontre que $P$ divise $P_n$.
\end{proof}

\begin{thm}
Le polynôme $P_n(X) = X^{p^n} - X$ est sans facteur carré dans $\F_p[X]$ et on a la décomposition en facteurs irréductibles :
\[ X^{p^n} - X = \prod_{d \in \D_n} \prod_{P \in \U_d(p)} P . \]
\end{thm}

\begin{proof}
Par dérivation on a $P_n'(X) = p^n X^{p^n-1} - 1 = -1$ et donc $P_n$ et $P'_n$ sont premiers entre eux, ce qui démontre que $P_n$
est sans facteur carré. Du lemme $\ref{lem:PdivPn}$ on déduit la forme de la décomposition en facteurs irréductibles.
\end{proof}

\begin{coro}
On a la formule :
\[ p^n = \sum_{d \in \D_n} d \cdot I_d(p) . \]
\end{coro}

\begin{thm}
Pour tout entier naturel non nul $n$, on a :
\[ n I_n(p) = \sum_{d \in \D_n} \mu \left ( \dfrac{n}{d} \right ) p^d . \]
\end{thm}

\begin{proof}
On utilise la formule d'inversion de Möbius :
\[ p^n = \sum_{d \in \D_n} d \cdot I_d(p) \iff n I_n(p) = \sum_{d \in \D_n} \mu(d) p^{n/d} 
	= \sum_{d \in \D_n} \mu \left ( \dfrac{n}{d} \right ) p^d , \]
la dernière égalité étant obtenu en considérant que $d \mapsto n/d$ est une permutation des éléments de $\D_n$.
\end{proof}

\begin{coro}
Pour tout entier naturel, il existe dans $F_p[X]$ des polynômes irréductibles de degré $n$ et on a :
\[ I_n(p) \sim_{n \to \infty} \dfrac{p^n}{n} . \]
\end{coro}

\begin{proof}
Pour $n=1$ on a $I_1(p) = p \geq 2$ et
pour $n=2, I_2(p) = \dfrac{p^2 - p}{2} \geq 1$. \\
Pour $n \geq 3$, on a :
\[ n \cdot I_n(p) = p^n + \sum_{d \in \D_n \setminus \set{n}} \mu \left ( \dfrac{n}{d} \right ) p^d . \]
On remarque que pour tout $d \in \D_n \setminus \set{n}$ on a $n=qd$ avec $2 \leq q \leq n$ soit $1 \leq d \leq E(n/2)$.
De plus $\abs*{\mu \left ( \dfrac{n}{d} \right )} \leq 1$, donc :
\begin{align*}
\abs*{\sum_{d \in \D_n \setminus \set{n}} \mu \left ( \dfrac{n}{d} \right ) p^d} & \leq \sum_{k=1}^{E(n/2)} p^d \\
& \leq p \dfrac{p^{E(n/2)} - 1}{p-1} \\
& < p^{E(n/2)+1} \leq p^{n/2+1} .
\end{align*}
On déduit donc que $n I_n(p) > p^n - p{n/2+1} \geq 0$ et :
\[ \abs*{\dfrac{n \cdot I_n(p)}{p^n} - 1} < \dfrac{p^{n/2+1}}{p^n} = p^{1-n/2} \to_{n \to \infty} 0 , \]
soit l'équivalent $I_n(p) \sim_{n \to \infty} \dfrac{p^n}{n}$.
\end{proof}

\begin{rmq}
Ce dernier corollaire nous dit qu'on peut construire un corps comportement exactement $p^n$ éléments,
à savoir le corps $\F_{p^n} = \F_p[X] / (P)$, où $P \in U_n(p)$.
\end{rmq}

\begin{thm}
A un isomorphisme près, il n'existe qu'un seul corps à $p^n$ éléments, c'est le corps $\F_{p^n} = \F_p[X] / (P)$,
où $P \in U_n(p)$.
\end{thm}

\begin{proof}
On désigne par $\K$ un corps fini à $p^n$ éléments, donc de caractéristique $p$
Le sous-corps premier de $\K$ est donc isomoprhe à $\F_p$,
et on peut donc identifier tout polynôme dans $\F_p[X]$ à un polynôme dans $\K[X]$. \newline

Selon le lemme $\ref{lem:PdivPn}$, le polynôme $P$ divise $P_n(X) = X^{p^n} - X$ dans $F_p[X]$ donc dans $\K[X]$. 
Or, selon le théorème de Lagrange, comme $\card{\Ket} = p^n-1$, on a $\lambda^{p^n} = \lambda$ 
pour tout $\lambda$ dans $\K$ donc  le polynôme $P_n$ est nul pour tout élément de $\K$, on peut donc écrire :
\[ P_n(X) = \prod_{\lambda \in \K} (X - \lambda) . \]
$P_n$ est scindé et à racines simples dans $\K[X]$, il en est donc de même pour $P$ qui divise $P_n$. \newline

Considérons maintenant l'application :
\[ \begin{array}{c c c c}
f : & \F_p[X] & \to & \K \\
& Q & \mapsto & Q(\lambda) , 
\end{array} \]
avec $\lambda \in \K$ racine de $P$. $f$ est un morphisme d'anneaux et son noyau est donc un idéal de $\K[X]$.
Comme l'anneau $\K[X]$ est principal, il existe $Q \in \F_p[X]$ unitaire et tel que $\ker{f} = (Q)$.
De plus, l'idéal $(P)$ est inclus dans $(Q)$, donc le polynôme $Q$ divise le polynôme unitaire irréductible $P$ dans $\F_p[X]$.
Le cas $Q=1$ est exclu, car alors $\lambda$ est racine de tout polynôme de $\F_p[X]$ ce qui est absurde. 
On conclut donc que $P=Q$ et, par suite, $\ker{f} = (P)$. La restriction de $f$ à $\F_{p^n}$ est donc une injection dans $\K$.
Ces deux corps étant de même cardinal, on conclut qu'ils sont isomorphes.
\end{proof}


\end{document}